\documentclass[11pt]{article}
\usepackage[a4paper,margin=1in]{geometry}
\usepackage[T1]{fontenc}
\usepackage[utf8]{inputenc}
\usepackage{lmodern}
\usepackage{enumitem}
\usepackage{hyperref}
\usepackage{setspace}
\setstretch{1.08}
\setlength{\parindent}{0pt}
\setlength{\parskip}{0.55em}

\begin{document}

\begin{center}
{\LARGE Response Letter}
\end{center}

We thank the Editor and the Reviewers for the careful reading and constructive feedback.
We have revised the manuscript accordingly and clarified the scope of our claims.
Below we provide a point-by-point response.

\section*{Reviewer 1}

\subsection*{Comment 1}
\begin{quote}
The method largely reuses existing multi-label supervised contrastive learning with metadata labels.
The current framing sometimes reads as if the loss is novel rather than the application choice and evaluation.
\end{quote}

\textbf{Response.}
Thank you for this important point. We agree and have revised the manuscript to avoid implying that the loss itself is novel.
Our contribution is now consistently framed as:
\begin{itemize}[leftmargin=1.5em,noitemsep,topsep=0.3em]
\item a metadata-supervised pretraining application using established multi-label supervised contrastive learning, and
\item a systematic transfer evaluation on three downstream tasks.
\end{itemize}
We updated wording in the Abstract, Introduction, Experimental Setting, and table captions to make this distinction explicit.

\subsection*{Comment 2}
\begin{quote}
The study evaluates ModAn-MulSupCon only with a ResNet-18 and a 2D setting.
The authors also acknowledge this limitation in the Discussion, but no evidence is provided on whether the reported under stronger architectures.
\end{quote}

\textbf{Response.}
Thank you for raising this important limitation.
We fully agree that evidence from stronger architectures would further strengthen the study.
In this revision cycle, however, we were not able to add new large-scale experiments with stronger backbones because of resource and scope constraints.
Each pretraining run requires approximately 10 minutes per epoch (as stated in the manuscript), and we currently have access to only a single GPU, which limited the number of additional experiments we could complete.

To avoid over-interpretation, we revised the manuscript to:
\begin{itemize}[leftmargin=1.5em,noitemsep,topsep=0.3em]
\item explicitly limit claims to the tested setting (ResNet-18, 2D, miniRIN pretraining), and
\item clearly state this as a limitation in Discussion.
\end{itemize}
We also added concrete future-work items: evaluation with ViT-style backbones and 3D encoders, and broader multi-institutional settings.

\subsection*{Comment 3}
\begin{quote}
Some recent medical image work should be discussed in the Related work, such as [1][2] doing pre-training on graph models or ConvNet, [2] doing pre-training on multimodal MRI data.
\begin{itemize}[leftmargin=1.5em,noitemsep,topsep=0.3em]
\item [1] Multi-modal hypergraph contrastive learning for medical image segmentation, PR
\item [2] Hypergraph Tversky-Aware Domain Incremental Learning for Brain Tumor Segmentation with Missing Modalities, MICCAI
\item [3] Cancer survival prediction from whole slide images with self-supervised learning and slide consistency, TMI
\end{itemize}
\end{quote}

\textbf{Response.}
Thank you for the helpful suggestions.
In this revision, we added the three suggested references to the Related Works section.
We also included a brief clarification of their relevance to our study.
At the same time, to avoid over-claiming, we did not add an extensive comparative discussion in this revision.
Our manuscript now states more modestly that these studies address related but different settings, while our work focuses on metadata-supervised contrastive pretraining for classification transfer.

\subsection*{Comment 4}
\begin{quote}
Each sample appears to have exactly two active labels (one modality, one anatomy).
This should be clearly explained.
\end{quote}

\textbf{Response.}
Thank you for this valuable suggestion.
In the revised manuscript, we added a dedicated subsection, \emph{Multi‑Hot Label Construction}, in the Method section to explicitly describe how the labels are defined.
We also added a concrete miniRIN example (ultrasound thyroid image) to illustrate this construction.

\subsection*{Comment 5}
\begin{quote}
The manuscript does not quantify how many positives per anchor arise, the mechanism of contrastive learning needs further clarification.
\end{quote}

\textbf{Response.}
Thank you for this key request.
In the revised manuscript, we report a concise summary of this point, while we provide the detailed analysis below as reference information in this response letter.
The concise summary has been added in Section 4.1.1 (Pretraining Dataset).
We quantified positives per anchor under the default setting (\(\tau=0.3\), batch size \(256\), two-view setup).

Main statistics (973,260 anchors):
\begin{itemize}[leftmargin=1.5em,noitemsep,topsep=0.3em]
\item positives per anchor: mean \(329.814\) (candidate pool \(2N-1=511\)),
\item positive ratio \(|P_\tau(a)|/(2N-1)\): mean \(0.6478\), median \(0.6908\),
\end{itemize}

Mechanism breakdown (mean extra positives per anchor):
\begin{itemize}[leftmargin=1.5em,noitemsep,topsep=0.3em]
\item modality-only matches: \(216.512\),
\item anatomy-only matches: \(53.179\),
\item both modality+anatomy matches: \(59.122\).
\end{itemize}

\subsection*{Comment 6}
\begin{quote}
In Eq. (1)-(3), the contrastive losses use the negatives-only denominator.
The manuscript does not explicitly state whether positive samples are included in the denominator.
\end{quote}

\textbf{Response.}
Thank you for identifying the ambiguity.
In our implementation, the denominator includes all non-self samples (positives and negatives), consistent with standard InfoNCE/SupCon.

\section*{Reviewer 2}

\subsection*{Comment 1}
\begin{quote}
Please describe more detail of your method in section 3.
\end{quote}

\textbf{Response.}
Thank you for this valuable suggestion.
To address this point, we added a dedicated subsection, \emph{Multi‑Hot Label Construction}, in Section 3 (Method).

\subsection*{Comment 2}
\begin{quote}
The size of the images of each is no more less than 1000, if possible, please add more dataset or increase more images to each dataset.
\end{quote}

\textbf{Response.}
Thank you for this suggestion.
We agree that larger downstream datasets would further strengthen the conclusions.

In this revision, we kept the three established public downstream tasks to preserve protocol comparability, and we already used repeated \(10\times5\)-fold CV to reduce variance.
We now state this sample-size constraint more explicitly as a limitation.

We also strengthened the future-work plan to include:
\begin{itemize}[leftmargin=1.5em,noitemsep,topsep=0.3em]
\item adding larger and more diverse downstream datasets,
\item scaling pretraining data beyond miniRIN,
\item broader validation across institutions and modalities.
\end{itemize}

\end{document}
